% 'nonacm' to remove any running header
\documentclass[conference]{IEEEtran}

\usepackage{cite}
\usepackage{amsmath,amsfonts}
\usepackage{algorithmic}
\usepackage{graphicx}
\usepackage{textcomp}
\usepackage{xcolor}
\usepackage{booktabs}
\usepackage{tagging}
\usepackage{url}

\usepackage{sc26repro}

\pagenumbering{gobble} % no page numbers

% Uncomment/comment the following usetag lines
% if you would like to get an explanation or an example.
% Don't forget to comment them for the final submission.
% \usetag{explanation}
% \usetag{example}

\begin{document}

\twocolumn[%
{\begin{center}
\Huge
Appendix: Artifact Description
/ Artifact Evaluation
\end{center}}
]

%%%%%%%%%%%%%%%%%%%%%%%%%%%%%%%%%%%%%%%%%%%%%%%%%%%%%
%  AD Appendix
%%%%%%%%%%%%%%%%%%%%%%%%%%%%%%%%%%%%%%%%%%%%%%%%%%%%%

\appendixAD

\section{Overview of Contributions and Artifacts}

\subsection{Paper's Main Contributions}

\begin{description}
\item[$C_1$] We present FastMPK, a high-performance sparse matrix power kernel (MPK) that exploits cache reuse across successive SpMVs. FastMPK models the matrix as a dependency graph, recursively partitions it into cache-resident subgraphs, and employs a dependency-preserving runtime scheduler. Evaluated on two x86 CPUs, one ARM CPU, and two NVIDIA GPUs across 110 large SuiteSparse matrices, FastMPK achieves average speedups of up to 1.84$\times$ over state-of-the-art cache-reuse MPK implementations and up to 1.43$\times$ over vendor-tuned libraries.

\item[$C_2$] We demonstrate the practical impact of FastMPK by integrating it as a polynomial preconditioner into the GMRES solver in Ginkgo, achieving end-to-end speedups of up to 1.68$\times$ on CPU and 1.40$\times$ on GPU with minimal preprocessing overhead.
\end{description}

\subsection{Computational Artifacts}

\begin{description}
\item[$A_1$] FastMPK --- CPU and GPU implementations \\
\url{https://github.com/qleonardo/FastMPK} \\
Zenodo: \url{https://doi.org/10.5281/zenodo.22089068}

\item[$A_2$] FastMPK-integrated GMRES --- CPU and GPU implementations \\
\url{https://github.com/qleonardo/FastMPK-GMRES} \\
Zenodo: \url{https://doi.org/10.5281/zenodo.22089070}

\end{description}

\begin{center}
\begin{tabular}{rll}
\toprule
Artifact ID  &  Contributions &  Related \\
             &  Supported     &  Paper Elements \\
\midrule
$A_1$   &  $C_1$ & Table III \\
        &        & Figures 2, 3, 4, 5, 6 \\
\midrule
$A_2$   &  $C_2$ & Figure 7 \\
\bottomrule
\end{tabular}
\end{center}


%%%%%%%%%%%%%%%%%%%%%%%%%%%%%%%%%%%%%%%%%%%%%%%%%%%%%%%%
\section{Artifact Identification}
%%%%%%%%%%%%%%%%%%%%%%%%%%%%%%%%%%%%%%%%%%%%%%%%%%%%%%%%

\newartifact
\artrel

$A_1$ (FastMPK) is the core computational artifact implementing the sparse matrix power kernel. It includes the following MPK variants for performance comparison:

\begin{itemize}
  \item \textbf{CSR}: baseline MPK using our CSR-based SpMV.
  \item \textbf{MKL} (CPU) / \textbf{cuSPARSE} (GPU): MPK using vendor-optimized SpMV from Intel oneMKL or NVIDIA cuSPARSE.
  \item \textbf{FBMPK} (CPU only): the cache-reuse MPK by Zhang et al.\ (IPDPS 2023). This is an external baseline; DOI: \url{https://doi.org/10.1109/IPDPS54959.2023.00046}. The source code is not publicly available; we obtained it directly from the original authors upon request and include it in our artifact repository with their permission.
  \item \textbf{RACE} (CPU only): the RACE-based cache-reuse MPK by Alappat et al.\ (IEEE TPDS 2022). External baseline; DOI: \url{https://doi.org/10.1109/TPDS.2022.3223512}. Code: \url{https://github.com/RRZE-HPC/RACE}; Zenodo: \url{https://doi.org/10.5281/zenodo.13071390}.
  \item \textbf{FastMPK (w/o Level)} (CPU only): FastMPK without the level-based traversal optimization. Not available on GPU, consistent with the paper.
  \item \textbf{FastMPK}: the full version with all optimizations enabled.
\end{itemize}

Both CPU and GPU implementations are included. The code is compiled in double precision.

\artexp

The measured performance is reported in GFLOPS. FastMPK should outperform all baselines on most matrices, with larger margins for higher powers $s$ and matrices with high reuse ratios. The level-based optimization should further improve performance over FastMPK (w/o Level). These results substantiate the cache-reuse strategy and dependency-aware parallelization presented in the paper.

\arttime

\begin{center}
\begin{tabular}{ll}
\toprule
Platform & Estimated Time (mins) \\
\midrule
Intel Xeon Max 9468 (CPU)  & $\sim$180 \\
AMD EPYC 7742 (CPU)        & $\sim$210 \\
Huawei Kunpeng 920 (CPU)   & $\sim$240 \\
NVIDIA H100 (GPU)          & $\sim$120 \\
NVIDIA A100 (GPU)          & $\sim$180 \\
\bottomrule
\end{tabular}
\end{center}
Times are for the 110-matrix benchmark with $s\in\{5,10\}$. Actual duration varies with the number of matrices, $s$ values, and available cores. The provided \texttt{run.sh} uses the 7 representative matrices for quick verification (estimated 10--15~mins).

\artin

\artinpart{Hardware}

\emph{CPU}: Intel Xeon Max 9468 @ 2.1\,GHz (48 cores, HBM-only mode), AMD EPYC 7742 @ 1.5\,GHz (64 cores), Huawei Kunpeng 920 @ 2.6\,GHz (32 cores).

\emph{GPU}: NVIDIA H100 SXM5 80\,GB and NVIDIA A100 PCIe 40\,GB.

\artinpart{Software}

\emph{CPU}: Intel oneAPI 2025.1 with oneMKL (Intel) or GCC $\ge$ 9 with the oneMKL runtime, OpenMP, Python 3 (for \texttt{download\_matrices.py}).

\emph{GPU}: CUDA 12.1 (H100) or CUDA 11.7 (A100), cuSPARSE.

\artinpart{Datasets / Inputs}

A Python script (\texttt{download\_matrices.py}) is provided at the repository root to download the 7 representative matrices from the SuiteSparse Matrix Collection\footnote{\url{https://sparse.tamu.edu/}}. Run \texttt{python download\_matrices.py} to populate the \texttt{matrix/} directory. The full evaluation uses 110 large sparse matrices (nnz $\ge$ 8,000,000) from the same collection. An input vector is randomly generated for each matrix. The matrix power $s$ ranges from 1 to 10 (key results at $s=5$ and $s=10$).

\artinpart{Installation and Deployment}

\emph{CPU}: Compile with \texttt{make} (default: \texttt{icpx} + oneMKL) or \texttt{make CC=g++}, which automatically switches to \texttt{-fopenmp} and the GNU-threaded MKL runtime. Set thread affinity with \texttt{OMP\_PROC\_BIND=close} and \texttt{numactl --physcpubind=... --membind=0}, one thread per physical core (no hyperthreading). On Xeon Max 9468, use HBM-only mode.

\emph{GPU}: Compile with \texttt{make} (default architecture \texttt{sm\_80} for A100; use \texttt{make ARCH=90} for H100).

\artcomp

For each matrix and each MPK variant, the program runs 100 iterations and reports the average execution time. A batch script (\texttt{run.sh}) is provided for quick verification: it automatically detects the thread count and NUMA topology, configures the binding accordingly (and simply skips binding with a warning if the topology cannot be detected), then invokes the executable over the test matrices. All machine-specific settings can be overridden via environment variables (\texttt{MATRIX\_DIR}, \texttt{OMP\_NUM\_THREADS}, \texttt{NUMA\_CMD}) without editing the script. For a single matrix, run:
\begin{center}
\texttt{./main <matrix.mtx> <s>}
\end{center}
The matrix name is automatically extracted from the file path.

\artout

Each run appends to \texttt{performance\_result.csv}: matrix name, dimensions, and GFLOPS for each variant. At the end of each matrix, the program directly prints a summary of FastMPK speedups over the baselines, e.g.:
\begin{center}
\texttt{FastMPK vs CSR: 2.10x ~ ~ FastMPK vs MKL: 2.11x}
\end{center}
The key metric is GFLOPS; FastMPK should achieve the highest values on most matrices.


\newartifact

\artrel

$A_2$ integrates FastMPK as a polynomial preconditioner into the GMRES solver from the Ginkgo (v1.11.0) library\footnote{\url{https://github.com/ginkgo-project/ginkgo}}. Two solver configurations are compared:

\begin{itemize}
  \item \textbf{Ginkgo (baseline)}: polynomial preconditioned GMRES using Ginkgo's standard CSR SpMV called sequentially.
  \item \textbf{FastMPK}: the same solver with polynomial evaluation replaced by FastMPK for cache reuse.
\end{itemize}

The preconditioner uses a Leja-ordered Newton basis with degree $d=40$. Both CPU (\texttt{OmpExecutor}) and GPU (\texttt{CudaExecutor}) implementations are provided. This artifact demonstrates that FastMPK can be easily integrated to accelerate real iterative solvers ($C_2$).

\artexp

Both configurations perform identical arithmetic, so they converge in the same number of iterations. FastMPK should reduce execution time due to lower memory traffic. Expected end-to-end speedups are 1.20$\times$--2.86$\times$ on CPU (geometric mean $\sim$1.68$\times$) and $\sim$1.40$\times$ on GPU. The one-time preprocessing cost should be under 5\% of total solver time.

\arttime

\begin{center}
\begin{tabular}{ll}
\toprule
Platform & Estimated Time (mins) \\
\midrule
Intel Xeon Max 9468 (CPU) & $\sim$20 \\
NVIDIA H100 (GPU)         & $\sim$10 \\
\bottomrule
\end{tabular}
\end{center}

\artin

\artinpart{Hardware}

\emph{CPU}: Intel Xeon Max 9468 @ 2.1\,GHz (48 cores, HBM-only mode).

\emph{GPU}: NVIDIA H100 SXM5 80\,GB.

\artinpart{Software}

Ginkgo v1.11.0 (\url{https://github.com/ginkgo-project/ginkgo}), CMake $\ge$ 3.16, OpenMP, Python 3 (for \texttt{download\_matrices.py}). CPU: Intel oneAPI 2025.1 or GCC $\ge$ 9. GPU: CUDA 12.1 with cuBLAS and cuSPARSE.

\artinpart{Datasets / Inputs}

The 7 representative matrices listed in Table~II of the paper (af\_shell10, audikw\_1, Geo\_1438, gsm\_106857, Serena, Long\_Coup\_dt0, SiO2). Use \texttt{python download\_matrices.py} to download them into the \texttt{matrix/} directory. A random RHS vector $\mathbf{b}$ is generated for each.

\artinpart{Installation and Deployment}

Both CPU and GPU versions use CMake in Release mode (\texttt{-O3}):
\begin{center}
\texttt{cmake -B build -DGinkgo\_DIR=<path> \&\& cmake --build build}
\end{center}
On GPU, additionally set \texttt{-DCMAKE\_CUDA\_ARCHITECTURES=90} for H100 (or 80 for A100).

\artcomp

For each matrix, the program reads the \texttt{.mtx} file, computes $\lambda_{\max}$ via power iteration, and runs each GMRES configuration 100 times. A batch script (\texttt{run.sh}) is provided for quick verification: it automatically detects the thread count and NUMA topology, configures the binding accordingly (and simply skips binding with a warning if the topology cannot be detected), then invokes the solver over the test matrices. All machine-specific settings can be overridden via environment variables (\texttt{MATRIX\_DIR}, \texttt{OMP\_NUM\_THREADS}, \texttt{NUMA\_CMD}) without editing the script. The solver uses Krylov dimension 100, relative residual tolerance $10^{-6}$, maximum 1,000 iterations, and polynomial degree $d=40$.

\artout

The program outputs to \texttt{stdout}: variant name, iteration count, execution time, and the end-to-end speedup of FastMPK over Ginkgo (printed directly at the end, e.g. \texttt{FastMPK end-to-end speedup over Ginkgo: 2.20x}). Both variants should report the same iteration count.


%%%%%%%%%%%%%%%%%%%%%%%%%%%%%%%%%%%%%%%%%%%%%%%%%%%%%
%  AE Appendix
%%%%%%%%%%%%%%%%%%%%%%%%%%%%%%%%%%%%%%%%%%%%%%%%%%%%%
\newpage
\appendixAE

\arteval{1}
\artin

\artinpart{Obtaining the Artifact}

Download from: \url{https://github.com/qleonardo/FastMPK}, or use the Zenodo archive at \url{https://doi.org/10.5281/zenodo.22089068}.

\artinpart{Software Dependencies}

\emph{CPU}: Intel oneAPI 2025.1 with MKL (Intel), or GCC $\ge$ 11.2 (AMD/ARM), with OpenMP.

\emph{GPU}: CUDA Toolkit $\ge$ 11.7 with cuSPARSE.

\artinpart{Datasets}

Run \texttt{python download\_matrices.py} to fetch the 7 representative matrices from the SuiteSparse Matrix Collection\footnote{\url{https://sparse.tamu.edu/}} in Matrix Market (\texttt{.mtx}) format into the \texttt{matrix/} directory.

\artinpart{Building the Artifact}

\emph{CPU}: Run \texttt{make} (default: Intel oneAPI \texttt{icpx} + MKL), or \texttt{make CC=g++} on GCC-only systems, which automatically switches to \texttt{-fopenmp} and the GNU-threaded MKL runtime.

\emph{GPU}: Run \texttt{make} (default architecture \texttt{sm\_80} for A100; use \texttt{make ARCH=90} for H100).

All platforms produce a single executable \texttt{main} (double precision).

\artinpart{Deployment}

The \texttt{run.sh} script handles deployment automatically: it detects the thread count and NUMA topology, configures the binding accordingly, and skips binding with a warning if the topology cannot be detected. All settings can be overridden via environment variables (\texttt{MATRIX\_DIR}, \texttt{OMP\_NUM\_THREADS}, \texttt{NUMA\_CMD="none"} to disable binding) without editing the script. On Xeon Max 9468, configure HBM-only mode.

\artcomp

Run \texttt{bash run.sh <s>} (or submit it via \texttt{sbatch}); the script automatically configures threads and NUMA binding and invokes the executable over the test matrices. For a single matrix (CPU):
\begin{center}
\texttt{OMP\_NUM\_THREADS=<N> OMP\_PROC\_BIND=close numactl --physcpubind=... --membind=0 ./main <matrix.mtx> <s>}
\end{center}
For GPU:
\begin{center}
\texttt{./main <matrix.mtx> <s>}
\end{center}
where \texttt{<s>} is the matrix power (typically 5 or 10). The matrix name is extracted from the file path. The program runs each MPK variant 100 times and reports average GFLOPS.

\artout

Results are appended to \texttt{performance\_result.csv} with columns: matrix name, rows, nnz, and GFLOPS for each MPK variant. At the end of each matrix, the program prints a summary of FastMPK speedups over the baselines (CSR, MKL/cuSPARSE, and FastMPK w/o Level on CPU). FastMPK should achieve the highest GFLOPS on most matrices.


\arteval{2}
\artin

\artinpart{Obtaining the Artifact}

Download from: \url{https://github.com/qleonardo/FastMPK-GMRES}, or use the Zenodo archive at \url{https://doi.org/10.5281/zenodo.22089070}.

\artinpart{Software Dependencies}

CMake $\ge$ 3.16, Ginkgo v1.11.0 (\url{https://github.com/ginkgo-project/ginkgo}), OpenMP. CPU: Intel oneAPI 2025.1 or GCC $\ge$ 11.2. GPU: CUDA $\ge$ 11.7 with cuBLAS and cuSPARSE.

\artinpart{Datasets}

Download the 7 matrices listed in Table~II of the paper from SuiteSparse, in \texttt{.mtx} format. Use \texttt{python download\_matrices.py} to fetch them into the \texttt{matrix/} directory.

\artinpart{Building the Artifact}

From the artifact root directory:
\begin{center}
\texttt{cmake -B build -DGinkgo\_DIR=<path> \&\& cmake --build build}
\end{center}
where \texttt{<path>} points to the Ginkgo CMake configuration (e.g., \texttt{/usr/local/lib/cmake/Ginkgo}). For GPU, add \texttt{-DCMAKE\_CUDA\_ARCHITECTURES=90} (H100) or \texttt{=80} (A100). The build produces \texttt{build/bench}.

\artinpart{Deployment}

The \texttt{run.sh} script handles deployment automatically: it detects the thread count and NUMA topology, configures the binding accordingly, and skips binding with a warning if the topology cannot be detected. All settings can be overridden via environment variables (\texttt{MATRIX\_DIR}, \texttt{OMP\_NUM\_THREADS}, \texttt{NUMA\_CMD="none"} to disable binding) without editing the script.

\artcomp

Run \texttt{bash run.sh}; the script automatically configures threads and NUMA binding and invokes the solver over the test matrices. For a single run (CPU):
\begin{center}
\texttt{OMP\_NUM\_THREADS=<N> OMP\_PROC\_BIND=close numactl --physcpubind=... --membind=0 ./build/bench <matrix.mtx>}
\end{center}
For GPU:
\begin{center}
\texttt{./build/bench <matrix.mtx>}
\end{center}
The program reads the matrix, computes $\lambda_{\max}$, runs both Ginkgo and FastMPK GMRES configurations 100 times each, and reports average times. Solver parameters are hardcoded: $d=40$, Krylov dimension 100, tolerance $10^{-6}$, max 1,000 iterations.

\artout

Output to \texttt{stdout}: variant name, iteration count, execution time in seconds, and the end-to-end speedup of FastMPK over Ginkgo (printed directly at the end). Both variants must report identical iteration counts. Expected speedups: 1.20$\times$--2.86$\times$ on CPU, $\sim$1.40$\times$ geometric mean on GPU. The preprocessing time (printed separately) should be under 5\% of total solver time.

\end{document}
